% ============================================================
% 片段：s03a  对应原文 PDF 第 12--18 页
% 内容：§3 引言、3.1--3.4
% 译者：t-s03a（DeepSeek Harness 智能体）   审校状态：初稿
% ============================================================

% ==== 原文 PDF 第 12 页 ====

\section{拓扑弦理论与镜像对称}
\label{sec:3}

\begin{keytakeaways}
  \begin{itemize}[leftmargin=1.4em,itemsep=2pt,topsep=2pt]
    \item \textbf{物理图像}：拓扑弦理论是二维拓扑场论与二维引力的耦合。固定世界面上的拓扑场论没有动力学自由度；耦合引力后必须对一切亏格的黎曼面及其模空间求和，配分函数按亏格展开，其中 $g=0$ 的振幅即预势。
    \item \textbf{核心机制}：拓扑扭转把 $U(1)$ 流混入能动张量，把费米性的超流荷改造成格拉斯曼标量 BRST 荷，使 BRST 荷能在任意黎曼面上整体定义。这是由 $\mathcal{N}=2$ 超共形场论构造拓扑弦的关键一步。
    \item \textbf{两条路线}：限制到 $(a,c)$ 手征环得到 A 模型，其形变是复化凯勒类的形变；限制到 $(c,c)$ 手征环得到 B 模型，其形变是复结构的形变。镜像对称（$J\to-J$）交换这两条路线。
    \item \textbf{A 模型的结论}：$Q_A$ 的上同调等同于 de Rham 上同调，作用量可拆成 $Q_A$-恰当项与拓扑项；路径积分局域化到全纯映射，给出 Gromov--Witten 不变量，并使经典相交环被瞬子修正为量子上同调环。
    \item \textbf{数学结构}：形变子环上的联络 $\nabla_A$ 与三点函数（汤川耦合）的“经典三重相交 + 瞬子求和”分解，是镜像对称计算的核心对象。
  \end{itemize}
\end{keytakeaways}

超共形代数（superconformal algebra）的两个同构手征环（chiral ring）——$(a,c)$ 环与 $(c,c)$ 环——的几何实现，借助拓扑弦理论（topological string theory）来研究；后者是一个与引力耦合的二维拓扑场论（topological field theory）。

% [Tutor Note: 物理图像] 拓扑场论与拓扑弦理论的差别全部在“引力”上
\begin{yizhu}[物理图像]
拓扑场论与拓扑弦理论的差别全部在“引力”上：固定世界面上的拓扑场论没有动力学自由度，其关联函数只依赖插入算符的上同调类；一旦耦合二维引力，就必须对亏格 $g$ 的黎曼面及其 $3g-3$ 维模空间积分（见 (3.5)），从而得到按亏格组织的无穷族振幅。式 (3.5) 处的脚注特别指出 $g=0,1$ 需单独处理，其中 $\mathcal{F}^0$ 是预势（prepotential），它是镜像对称中计算汤川耦合的出发点。
\end{yizhu}

% [Tutor Note: 含糊点] 两个手征环“同构”不等于几何上等同
\begin{yizhu}[含糊点]
$(a,c)$ 与 $(c,c)$ 手征环作为手征代数同构，但它们在非线性 sigma 模型中的几何实现截然不同：前者对应复化凯勒形式，后者对应复结构。原文只用一句“借助拓扑弦理论研究”概括，真正的等价陈述（镜像对称）要到 3.3 节才出现；阅读时不要把这里的“同构”误当成几何上的等同。
\end{yizhu}

% [Tutor Note: 易错点] 花体 \mathcal{Q} 与具体 BRST 荷 Q_A, Q_B 的记号分层
\begin{yizhu}[易错点]
记号是分层的：本节（第 12 页）用花体 $\mathcal{Q}$ 表示抽象的上同调型 BRST 算符，而 3.2 节起用 $Q_A,Q_B$（以及 $Q^{\pm},\bar Q^{\pm}$）表示扭转过后的具体荷。$\mathcal{Q}$-恰当性（$\mathcal{Q}$-exact）是本节一切结论（度规无关、半经典精确）的唯一来源，务必与具体模型中的 $Q_A$ 区分。
\end{yizhu}

\subsection{拓扑场论}
\label{sec:3.1}

下面概述上同调型（cohomological type）拓扑 sigma 模型 \cite{ref35} 的一些特征。它们是 sigma 模型，即从 $\Sigma_g$ 映入某个靶空间 $X$ 的映射 $\phi:\Sigma_g\to X$ 的理论；所关心的是 $\Sigma_g$ 表示亏格 $g$ 的黎曼面的情形。这些映射在物理上由玻色子与费米子实现，二者由超对称相联系。上同调型拓扑场论的额外结构，在于存在一个格拉斯曼（Grassmann）的、标量的对称算符 $\mathcal{Q}$（有时称为 BRST 算符），它具有如下性质：
\begin{itemize}[leftmargin=1.4em,itemsep=1pt,topsep=2pt]
  \item $\mathcal{Q}$ 幂零，$\mathcal{Q}^2=0$。
  \item 能动张量是 $\mathcal{Q}$ 恰当的，
\end{itemize}
\[
T_{\mu\nu}=\{\mathcal{Q},G_{\mu\nu}\}\,.
\]

后一性质的直接推论是：关联函数不依赖于二维度规。为看出这一点，注意对度规的变分等价于在关联函数中插入能动张量；后者是恰当的，而 $\mathcal{Q}$ 又是理论的对称性，这就说明了度规无关性。此外，理论的物理态对应于算符 $\mathcal{Q}$ 的上同调类，并且用于计算路径积分的半经典近似是精确的。更多的背景以及这些性质的推导可参见 \cite{ref36,ref9}。

% [Tutor Note: 物理图像] Q^2=0 使物理态空间成为 Q 的上同调
\begin{yizhu}[物理图像]
$\mathcal{Q}^2=0$ 使 $\mathcal{Q}$ 的上同调成为物理态空间：$\mathcal{Q}$-闭态模去 $\mathcal{Q}$-恰当态。相应地，$\mathcal{Q}$-恰当算符的插入在关联函数中给出零（或全导数），这正是拓扑场论只需上同调数据的根本原因；本节末尾“物理态对应于上同调类”即指此。
\end{yizhu}

% [Tutor Note: 含糊点] 度规无关性论证中被省略的 Ward 恒等式
\begin{yizhu}[含糊点]
度规无关性的论证被压缩成一句：度规变分等于插入 $T_{\mu\nu}$，而 $T_{\mu\nu}=\{\mathcal{Q},G_{\mu\nu}\}$ 是恰当的，$\mathcal{Q}$ 又是对称性，故度规变分为零。被省略的是 Ward 恒等式这一步：$\langle\{\mathcal{Q},G_{\mu\nu}\}\cdots\rangle=0$ 要求 $\mathcal{Q}$ 不 anomalous（无边界、无拓扑障碍）。若插入的算符本身不是 $\mathcal{Q}$-闭的，该论证失效。
\end{yizhu}

% [Tutor Note: 易错点] “半经典近似精确”与“只算经典解”不是一回事
\begin{yizhu}[易错点]
“半经典近似精确”不等于“只算经典解”：它的含义是路径积分在 $\mathcal{Q}$-不动点附近的高斯近似给出精确结果，即局域化。本节的论证只用到 $\mathcal{Q}$-恰当性；具体的局域化方程要到 3.4.1 的 (3.14)--(3.16) 才写出（$\mathcal{D}=0$ 等价于 $\bar\partial\phi=0$）。
\end{yizhu}

% ==== 原文 PDF 第 13 页 ====

\subsection{拓扑扭转与拓扑弦理论}
\label{sec:3.2}

拓扑场论可由 $\mathcal{N}=2$ 超共形场论构造，这类理论的态正是由手征环中的算符所产生的态。做法是把所论态的湮灭算符当作 BRST 算符，并且只把被 BRST 算符湮灭的态视为物理态。由式 (2.13) 中手征态的定义及其反手征对应物可以清楚看出：若要限制到 $(a,c)/(c,c)$ 手征环，则所论算符应为
\begin{equation}\tag{3.1}
\underset{(a,c)}{Q_A=G^{-}_{-1/2}+\bar G^{+}_{-1/2}}\,,\qquad
\underset{(c,c)}{Q_B=G^{+}_{-1/2}+\bar G^{+}_{-1/2}}\,.
\end{equation}

这些算符在代数中平方为零，$(Q_{A/B})^2=0$，从而满足 (3.1) 的第一条要求。% [原文疑误: 此处 “(3.1)” 应为 §3.1；$\mathcal{Q}^2=0$ 这一“第一条要求”是在 §3.1 的条目中提出的，(3.1) 本身只是 $Q_A,Q_B$ 的定义式，无“要求”可满足。照抄原文编号。]
然而，用其中任一算符写出的能动张量都不是 $\mathcal{Q}$ 恰当的；此外，这些算符是费米性的。

把所得理论耦合到二维引力即得到拓扑弦理论；这涉及对所有二维几何与度规的积分。二维几何即黎曼面，它们按亏格 $g$ 归类。特别地，必须能够在每一个这样的黎曼面上整体地定义 BRST 算符。然而，对应于 SCFT 的 $G^{\pm}$ 流的荷是费米子。为了在每一个黎曼面上都有整体良定义的荷，引入拓扑扭转（topological twist）。由于一个流的自旋由其共形权决定，而共形权出现在与能动张量的算符乘积展开中，扭转的想法便是修改能动张量。新的能动张量是旧能动张量与 $U(1)$ 流 $J$ 的组合。视所要限制的环而定，这通过下式实现：
\begin{equation}\tag{3.2}
\begin{array}{c|c}
A & B\\ \hline
T\to T+\frac{1}{2}\partial J & T\to T-\frac{1}{2}\partial J\\[6pt]
\bar T\to\bar T-\frac{1}{2}\bar\partial\bar J & \bar T\to\bar T-\frac{1}{2}\bar\partial\bar J\\ \hline
{\scriptstyle(a,c)} & {\scriptstyle(c,c)}
\end{array}
\end{equation}
从而得到格拉斯曼值的标量 $Q_A$ 与 $Q_B$。例如，限制到手征环时，现在有
\begin{equation}\tag{3.3}
T_{\mathrm{top}}=T-\frac{1}{2}\partial J\,,
\end{equation}
由 OPE (2.1) 可以看出，$G^+$ 现在的共形权为 $1$。此外还有
\begin{equation}\tag{3.4}
T_{\mathrm{top}}=\{Q^{+},G^{-}\}\,,
\end{equation}

% [Tutor Note: 物理图像] 引入扭转的动机是整体性
\begin{yizhu}[物理图像]
引入扭转的动机是整体性：$G^{\pm}$ 及其荷是时空费米子（自旋 $1/2$），在一般黎曼面上无法整体定义；把 $U(1)$ 流 $J$ 混入能动张量后，$G^{\pm}$ 的自旋改变，其荷变成格拉斯曼标量，才能作为整体定义的 BRST 算符。扭转不改变物理内容，只是重新选择世界面丛上的联络。
\end{yizhu}

% [Tutor Note: 易错点] (3.1) 中两个 BRST 荷的配对极易记错
\begin{yizhu}[易错点]
(3.1) 的配对极易记错：$(a,c)$ 环用 $Q_A=G^{-}_{-1/2}+\bar G^{+}_{-1/2}$，$(c,c)$ 环用 $Q_B=G^{+}_{-1/2}+\bar G^{+}_{-1/2}$；两式的反全纯部分都是 $\bar G^{+}_{-1/2}$，差别只在全纯部分是 $G^{-}$ 还是 $G^{+}$。若把 $Q_B$ 写成 $G^{+}_{-1/2}+\bar G^{-}_{-1/2}$ 就错了。
\end{yizhu}

% [Tutor Note: 易错点] (3.2) 表格的列与行方向
\begin{yizhu}[易错点]
(3.2) 表格的读法：左列 $A$ 对应 $(a,c)$、右列 $B$ 对应 $(c,c)$；两列的全纯扭转方向相反（$T\to T+\tfrac{1}{2}\partial J$ 与 $T\to T-\tfrac{1}{2}\partial J$），而反全纯部分两列相同，都是 $\bar T\to\bar T-\tfrac{1}{2}\bar\partial\bar J$。注意 $\partial$ 只作用于 $J$，$\bar\partial$ 只作用于 $\bar J$，第二行的两个横杠不可漏。
\end{yizhu}

% [Tutor Note: 原文疑误] “the first requirement of (3.1)”的编号疑为 §3.1 之误
\begin{yizhu}[原文疑误]
正文“满足 (3.1) 的第一条要求”中的 (3.1) 疑为 §3.1 之误：$\mathcal{Q}^2=0$ 这一“第一条要求”是在 §3.1 的列表中提出的，(3.1) 本身只是 $Q_A,Q_B$ 的定义式，并无“要求”可供满足。译文照录原文编号，供审校核对。
\end{yizhu}

% ==== 原文 PDF 第 14 页 ====

前提是关联函数中的 $U(1)$ 荷为 $c/3$，其中 $c$ 是中心荷。当超共形代数以映入卡拉比--丘（Calabi--Yau，CY）三维流形靶空间的非线性 sigma 模型几何实现时，$c=9$。这些算符的上同调给出到有限环的所需限制。正如记号所暗示的：限制到 $(a,c)$ 环给出 A 模型，限制到 $(c,c)$ 环给出 B 模型。

% [Tutor Note: 易错点] c/3 与 c=9 是 (3.3)(3.4) 成立的前提
\begin{yizhu}[易错点]
本段首句“关联函数中的 $U(1)$ 荷为 $c/3$”是 (3.3) (3.4) 成立的前提：CY 三维流形靶空间的非线性 sigma 模型有 $c=9$，故关联函数非零要求总 $U(1)$ 荷为 $c/3=3$。这个 $3$ 正是 (3.19) 三重相交项与手征环三点函数荷的来历；核对 §2 中 $J=\frac{1}{4}g_{i\bar\jmath}\psi^i\psi^{\bar\jmath}$ 的归一化后，$c/3$ 与算符荷的对应才能对上。
\end{yizhu}

与引力的耦合通过如下方式实现：对黎曼面的每个亏格定义关联函数，并把该关联函数在黎曼面模空间（moduli space）上积分。亏格 $g$ 黎曼面的模空间由 $3g-3$ 个 Beltrami 微分 $\mu^{z}_{\bar z}$ 描述；它们是取值于全纯切丛的反全纯形式，$\mu_a\in H^{0,1}(\Sigma_g,T\Sigma_g)$，以及它们的复共轭 $\bar\mu^{\bar z}_{z}$。这些微分与两个流缩并，后者的共形权被提升为 $h=2$（$\bar h=2$）。这一构造的动机来自玻色弦：在玻色弦中，反鬼场以类似的方式与 Beltrami 微分缩并。亏格 $g$ 的振幅或自由能现定义为（此处为 B 模型且 $g>1$）\footnote{$g=0,1$ 的振幅需要单独处理；特别地，$\mathcal{F}^0$ 表示预势，并可由形变丛的几何计算得到。}
\begin{equation}\tag{3.5}
\mathcal{F}^{g}=\int_{\mathcal{M}_g}[dm\,d\bar m]\left\langle
\prod_{a=1}^{3g-3}\left(\int_{\Sigma}\mu_a G^{-}\right)\left(\int_{\Sigma}\mu_{\bar a}\overline{G^{-}}\right)
\right\rangle_{\Sigma_g},
\end{equation}
其中 $\langle\cdots\rangle_{\Sigma_g}$ 表示 CFT 关联函数，$dm\,d\bar m$ 与 Beltrami 微分对偶。\footnote{关于拓扑弦的更多细节可参见 \cite{ref9,ref37,ref38,ref39,ref40,ref36,ref41}。}

\subsection{镜像对称}
\label{sec:3.3}

镜像对称（mirror symmetry）是把以 $(a,c)$ 与 $(c,c)$ 手征环为物理态的两个拓扑场（弦）理论等同起来。它在超共形代数中的起源是左移 $U(1)$ 流的变号：$J\leftrightarrow-J$。镜像对称的力量来自这样一个事实：在非线性 sigma 模型实现中，形变与手征环的几何解释非常不同。A 模型所看到的 $(a,c)$ 手征环的形变，是 CY 三维流形 $Z$ 的复化凯勒形式（complexified Kähler form）的形变。理论模空间上的物理坐标是\footnote{这里假设 $h^{2,0}(Z)=0$。}
\begin{equation}\tag{3.6}
t^{a}=\int_{C^{a}}B+i\mathcal{J}\,,\qquad a=1,\ldots,h^{1,1}(Z)\,,
\end{equation}
其中 $C^{a}\in H_2(Z)$，$B\in H^2(Z,\mathbb{C})/H^2(Z,\mathbb{Z})$，而 $\mathcal{J}$ 是凯勒形式（Kähler form）。

B 侧的形变是 CY 流形 $Z^{*}$ 的复结构（complex structure）形变。好的物理坐标是唯一全纯 $(3,0)$ 形式 $\tilde\Omega$ 的某些周期，并带有一定的归一化
\begin{equation}\tag{3.7}
t^{a}=\int_{\alpha^{a}}\tilde\Omega\,,\qquad a=1,\ldots,h^{2,1}(Z^{*})\,,
\end{equation}

% [Tutor Note: 物理图像] 镜像对称在超共形代数层面的操作
\begin{yizhu}[物理图像]
镜像对称在超共形代数层面的操作是左移 $U(1)$ 流变号 $J\to-J$，它把 $(a,c)$ 与 $(c,c)$ 手征环互换，从而把 A 模型的复化凯勒形变与 B 模型的复结构形变互换。几何上这给出两个不同 CY 流形 $Z$ 与 $Z^{*}$ 的模空间之间的映射（镜像映射），使一边的瞬子计算可以换成另一边的周期计算。
\end{yizhu}

% [Tutor Note: 易错点] (3.6) 中 B 场的大规范等价类与两侧同名的 t^a
\begin{yizhu}[易错点]
(3.6) 中 $t^{a}=\int_{C^{a}}(B+i\mathcal{J})$ 里的 $B$ 属于 $H^2(Z,\mathbb{C})/H^2(Z,\mathbb{Z})$：这是 $B$ 场的大规范等价类，不是普通 2-形式；$C^{a}\in H_2(Z)$ 是 2-循环基。(3.7) 中 $\alpha^{a}\in H_3(Z^{*})$ 是 3-循环，$\tilde\Omega$ 是 $Z^{*}$ 上唯一的 $(3,0)$ 形式。两侧都用 $t^{a}$ 是原文的记号，不要误以为二者定义在同一流形上。
\end{yizhu}

% [Tutor Note: 含糊点] 维数匹配只是镜像族的必要条件
\begin{yizhu}[含糊点]
$h^{1,1}(Z)=h^{2,1}(Z^{*})$ 与“形变相互匹配”只是镜像族的必要条件；真正的对应要求手征环（作为代数）同构，并且 $\tilde\Omega$ 的归一化要与手征环的几何实现一致。原文明确把归一化的确定推迟到后文，此处不能把 (3.7) 当作已经固定的约定。
\end{yizhu}

% ==== 原文 PDF 第 15 页 ====

其中 $\alpha^{a}\in H_3(Z^{*})$。复结构模空间上坐标的恰当选取以及 $\tilde\Omega$ 的归一化，由手征环的几何实现决定，并将在本讲义后面讨论。镜像对称所做的是找出 CY 空间的镜像族 $Z_t$ 与 $Z^{*}_{t^{*}}$，使得对族中每一成员都有 $h^{1,1}(Z)=h^{2,1}(Z^{*})$，且形变相互匹配。两侧模空间上正确物理坐标的对应，通过分别实现手征环的形变子环在两侧的几何对应物来达成。

\subsection{A 模型}
\label{sec:3.4}

\subsubsection{A 扭转}
\label{sec:3.4.1}

扭转 (3.2) 对靶空间为 $Z$ 的非线性 sigma 模型的作用，是改变费米子作为其截面的丛（参见 \cite{ref6,ref9,ref10}）
\begin{equation}\tag{3.8}
\begin{array}{ll}
\psi^{i}\in\Gamma(\phi^{*}T^{1,0}Z)\,, & \psi^{\bar\imath}_{z}\in\Gamma(K\otimes\phi^{*}T^{0,1}Z)\,,\\[5pt]
\chi^{i}_{\bar z}\in\Gamma(\bar K\otimes\phi^{*}T^{1,0}Z)\,, & \chi^{\bar\imath}\in\Gamma(\phi^{*}T^{0,1}Z)\,.
\end{array}
\end{equation}
其效果是：$\psi^{i}$ 与 $\chi^{\bar\imath}$ 成为非线性 sigma 模型二维底流形上的标量。为分析
\[
Q_A=Q^{-}+\bar Q^{+}
\]
的上同调，令 (2.32-2.37) 中 $\epsilon^{-}=\bar\epsilon^{+}=\epsilon$ 且 $\bar\epsilon^{-}=\epsilon^{+}=0$，于是超对称变分 $\delta$ 由
\[
\delta=\epsilon\,Q_A
\]
给出。标量量的变分变为（见 (2.32,2.34) 与 (2.33,2.37)）
\begin{equation}\tag{3.9}
\delta\phi^{i}=i\epsilon\,\psi^{i}\,,\qquad \delta\psi^{i}=0\,,
\end{equation}
\begin{equation}\tag{3.10}
\delta\phi^{\bar\imath}=i\epsilon\,\chi^{\bar\imath}\,,\qquad \delta\chi^{\bar\imath}=0\,,
\end{equation}
这提示如下对应：
\begin{equation}\tag{3.11}
\psi^{i}\leftrightarrow du^{i}\,,\qquad
\chi^{\bar\imath}\leftrightarrow d\bar u^{\bar\imath}\,,\qquad
Q_A\leftrightarrow d=\partial+\bar\partial\,,
\end{equation}
其中 $u^{i},\bar u^{\bar\imath}$ 是 $Z$ 上的局部坐标，$d$ 是 de Rham 算符。拓扑场论的算符与流形上的点相联系，因此只能由二维标量构成；这类算符应具有形式
\begin{equation}\tag{3.12}
\omega(\phi)_{i_1,\ldots,i_p,\bar\imath_1,\ldots,\bar\imath_q}\,
\psi^{i_1}\cdots\psi^{i_p}\,\chi^{\bar\imath_1}\cdots\chi^{\bar\imath_q}\,,
\end{equation}
它们可进而等同于
\begin{equation}\tag{3.13}
\omega(u)_{i_1,\ldots,i_p,\bar\imath_1,\ldots,\bar\imath_q}\,
du^{i_1}\cdots du^{i_p}\,d\bar u^{\bar\imath_1}\cdots d\bar u^{\bar\imath_q}\,.
\end{equation}

% [Tutor Note: 易错点] (3.8) 中四个截面的丛与指标配对
\begin{yizhu}[易错点]
(3.8) 的四个截面按 $(z,\bar z)$ 与 $(i,\bar\imath)$ 配对：$\psi^{i}$ 与 $\chi^{\bar\imath}$ 不带世界面微分（标量），$\psi^{\bar\imath}_{z}\in\Gamma(K\otimes\phi^{*}T^{0,1}Z)$、$\chi^{i}_{\bar z}\in\Gamma(\bar K\otimes\phi^{*}T^{1,0}Z)$ 分别带 $K$ 与 $\bar K$。原文用无点字母 $\imath,\jmath$ 表示反全纯指标（如 (3.15) 的 $g_{i\bar\jmath}$），转写时不能写成普通 $i,j$。
\end{yizhu}

% [Tutor Note: 物理图像] A 模型局域化的逻辑链
\begin{yizhu}[物理图像]
A 模型局域化的逻辑链是：作用量 (3.14) 的第一项是 $Q_A$-恰当的，把它的系数放大（$S\to S+t\{Q_A,\mathcal{D}\}$，$t\to\infty$）不改变路径积分，却使 $\mathrm{e}^{-t\{Q_A,\mathcal{D}\}}$ 只在 $\mathcal{D}=0$ 附近有贡献；由 (3.15)，$\mathcal{D}=0$ 即 $\bar\partial\phi=0$（全纯映射）。剩余积分是全纯映射模空间上的积分，其值即 Gromov--Witten 不变量；按靶空间 2-循环类 $\beta$ 分组就是对世界面瞬子扇区求和。
\end{yizhu}

% [Tutor Note: 含糊点] “these contributions are exact”背后的局域化定理
\begin{yizhu}[含糊点]
“it can be shown that these contributions are exact”指的是超对称局域化定理：$Q_A$-恰当形变不改变路径积分，因此单瞬子（高斯）近似已经精确。原文把它归入 localization 并指向 \cite{ref9,ref10}；被省略的是该定理的适用条件（$Q_A$ 的紧性、无边界、模空间的紧化），这正是 Gromov--Witten 理论需要处理的技术点。
\end{yizhu}

% [Tutor Note: 易错点] (3.15) 中 D 的缩并方向
\begin{yizhu}[易错点]
(3.15) 的缩并方向需要核对：第一项 $\psi^{\bar\jmath}_{z}\bar\partial\phi^{i}$ 中 $g_{i\bar\jmath}$ 把 $i$ 与 $\bar\jmath$ 配对，第二项 $\partial\phi^{\bar\jmath}\chi^{i}_{\bar z}$ 中 $i$ 与 $\bar\jmath$ 仍分别来自 $g_{i\bar\jmath}$ 的两个指标。若把 $\bar\partial$ 与 $\partial$（或 $z$ 与 $\bar z$）对调，$\mathcal{D}$ 就不再是 $Q_A$-恰当项的搭档，局域化方程也会改变。
\end{yizhu}

% ==== 原文 PDF 第 16 页 ====

拓扑理论的算符与 $Z$ 的 de Rham 上同调一一对应。此外，作用量现在可写成
\begin{equation}\tag{3.14}
S=i\int_{\Sigma}\{Q_A,\mathcal{D}\}-2\pi i\int_{\Sigma}\phi^{*}(B+i\mathcal{J})\,,
\end{equation}
其中凯勒形式为 $\mathcal{J}=ig_{i\bar\jmath}\,du^{i}\wedge d\bar u^{\bar\jmath}$，且
\begin{equation}\tag{3.15}
\mathcal{D}=2\pi g_{i\bar\jmath}\left(\psi^{\bar\jmath}_{z}\bar\partial\phi^{i}+\partial\phi^{\bar\jmath}\chi^{i}_{\bar z}\right).
\end{equation}
第二项只依赖于类 $\beta=\phi_{*}(\Sigma)\in H_2(Z,\mathbb{Z})$。因此配分函数与关联函数可写成
\begin{equation}\tag{3.16}
Z=\int[d\phi\,d\psi\,d\chi]\,\mathrm{e}^{-S}
=\sum_{\beta\in H_2}\mathrm{e}^{-2\pi i t_{\beta}}
\int_{\mathrm{fixed}\,\beta}[d\phi\,d\psi\,d\chi]\,
\mathrm{e}^{-i\int\{Q_A,\mathcal{D}\}}\,.
\end{equation}
通过把被积函数指数中的该项缩放至无穷大，可以论证主要贡献来自满足 $\mathcal{D}=0$ 的映射。这些是全纯映射
\[
\bar\partial\phi=0\,,
\]
并且可以证明这些贡献是精确的，这就是局域化（localization）的思想（更多细节参见 \cite{ref9,ref10} 及其中的文献）。在这些情形下，路径积分可定义为全纯映射模空间上的积分，并给出刻画靶空间 $Z$ 的数学不变量，即 Gromov--Witten 不变量。物理上，对绕靶空间中不同 2-循环的全纯映射求和，对应于对不同世界面瞬子（instanton）扇区求和。

\subsubsection{量子几何}
\label{sec:3.4.2}

前面已说明，由 A 扭转得到的拓扑场论的算符与 de Rham 上同调相对应，而且 $U(1)$ 荷被等同于形式的次数。手征环的几何解释由形式的外积给出，或等价地由 $Z$ 的相交环给出。然而，由于如式 (3.16) 那样对所有瞬子扇区求和，这个环被修正，给出量子上同调（quantum cohomology）环；它可看作相交环的量子力学修正，其中计入了全纯映射的计数。把量子力学的概念引入几何，称为量子几何（quantum geometry）。与经典几何理解的另一处偏离，是复化凯勒模空间的概念。经典上，凯勒形式的形变会定义凯勒锥。然而 sigma 模型表明，应当考虑的正确量是复化凯勒类以及扩大了的凯勒模空间；后者可按不同区域理解，这些区域称为相（phase），不同的 CY 几何可由这些相相联系 \cite{ref42,ref43,ref44}。

% [Tutor Note: 物理图像] 量子上同调 = 经典相交环 + 瞬子修正
\begin{yizhu}[物理图像]
量子上同调等于经典相交环加上瞬子修正。(3.19) 的第一项 $\int_Z\eta_a\wedge\eta_b\wedge\eta_c$ 是经典三重相交数（汤川耦合的经典值）；求和项对每个类 $\beta$ 以 $q^{[\beta]}/(1-q^{[\beta]})$ 加权，来自绕该有理曲线的所有多重覆盖（$\sum_{k\ge1}q^{k\beta}=q^{\beta}/(1-q^{\beta})$）。因此三点函数在 $q\to0$（大体积极限）回到经典值。
\end{yizhu}

% [Tutor Note: 易错点] (3.19) 的上标与下标约定
\begin{yizhu}[易错点]
(3.19) 的指标约定：$n_{[\beta]}$ 是下标（类 $\beta$ 的有理曲线数，由 (3.20) 组装成 Gromov--Witten 不变量 $N_{[\beta]}$），而 $q^{[\beta]}$ 是上标；$q^{[\beta]}=q_1^{c_1}\cdots q_n^{c_n}$，$c_a=\langle\beta,\eta_a\rangle$，$q_{\beta}:=\exp(2\pi i\int_{\beta}B+i\mathcal{J})$ 见 (3.19) 上方的无编号式。另需注意 (3.16) 中的权重写作 $\mathrm{e}^{-2\pi i t_{\beta}}$，与 $q_{\beta}$ 的定义相差一个符号约定，阅读时应按 $t_{\beta}$ 的定义核对。
\end{yizhu}

% [Tutor Note: 易错点] (3.18) 第二行的自由指标是 c
\begin{yizhu}[易错点]
(3.18) 第二行左边是 $\eta_c$ 而非 $\eta_a$：自由指标 $c$ 必须与 $C_{abc}$ 的第三个指标一致，右边对 $a,b$ 求和。若把左边写成 $\eta_a$，会造成指标重复。该联络是 Gauss--Manin 联络在形变子环上的实现，$q_a$ 是平坦坐标，$\nabla_A\eta_0=\sum_a\eta_a\otimes dq_a/q_a$ 即镜像映射的微分。
\end{yizhu}

% [Tutor Note: 含糊点] “large volume”展开的适用区域
\begin{yizhu}[含糊点]
“large volume”是瞬子求和 $\sum_{\beta}q^{[\beta]}$ 收敛的区域（$|q|<1$）。原文指出其他区域（如 orbifold 相）中镜像对称的数学表述仍在发展，因此 (3.19) 的 $q$ 展开不能直接解析延拓到整个模空间；相（phase）的讨论见本节前段与 \cite{ref42,ref43,ref44}。
\end{yizhu}

% ==== 原文 PDF 第 17 页 ====

下面简短阐述手征环的形变子环；更多的数学细节与背景可参见文献 \cite{ref45,ref8}。本阐述遵循 \cite{ref46}。量子上同调是在大体积（large volume）下展开的，这指的是 (3.16) 中瞬子求和具有意义的区域。模空间其他区域中镜像对称的数学表述仍在发展；例如关于涉及模空间中轨形（orbifold）区域的镜像对称，参见 \cite{ref47} 及其中文献。形变子环对应于 CY 流形 $Z$ 的偶上同调。复化凯勒类是 $H^2(Z,\mathbb{C})$ 中的元素。配对由形式的外积给出。为恢复手征环结构，应对下图给出解释
\begin{equation}\tag{3.17}
H^{0}(Z,\mathbb{C})\xrightarrow{\ \nabla_A\ }H^{2}(Z,\mathbb{C})
\xrightarrow{\ \nabla_A\ }H^{4}(Z,\mathbb{C})
\xrightarrow{\ \nabla_A\ }H^{6}(Z,\mathbb{C})\,.
\end{equation}
为定义该联络，令 $\eta_0$ 为 $H^0(Z,\mathbb{Z})$ 的生成元；$\eta_a$（$a=1,\ldots,n=h^{1,1}(Z)$）为 $H^2(Z,\mathbb{Z})$ 的一组基；$\chi_b$（$b=1,\ldots,n$）为 $H^4(Z,\mathbb{Z})$ 的一组基，它关于辛配对 $\langle\eta_a,\chi_b\rangle=\delta_{ab}$ 与前者对偶；最后令 $\chi_0\in H^6(Z,\mathbb{Z})$ 为 $\eta_0$ 的对偶。联络 $\nabla_A$ 现定义为
\begin{equation}\tag{3.18}
\begin{aligned}
\nabla_A\eta_0&=\sum_{a=1}^{n}\eta_a\otimes\frac{dq_a}{q_a}\,,\\
\nabla_A\eta_c&=\sum_{a,b=1}^{n}C_{abc}\,\chi_b\otimes\frac{dq_a}{q_a}\,,\\
\nabla_A\chi_b&=\chi_0\frac{dq_b}{q_b}\,,\\
\nabla_A\chi_0&=0\,.
\end{aligned}
\end{equation}
系数 $C_{abc}$ 是由有理曲线定义的 $q_1,\ldots,q_n$ 的幂级数，这些曲线的类由 $\beta\in H_2(Z,\mathbb{Z})$ 给出
\[
q_{\beta}:=\exp\left(2\pi i\int_{\beta}B+i\mathcal{J}\right),
\]
\begin{equation}\tag{3.19}
C_{abc}=\int_{Z}\eta_a\wedge\eta_b\wedge\eta_c
+\sum_{\beta\in H_2}n_{[\beta]}\langle\beta,\eta_a\rangle\langle\beta,\eta_b\rangle\langle\beta,\eta_c\rangle\,
\frac{q^{[\beta]}}{1-q^{[\beta]}}\,,
\end{equation}
其中 $q^{[\beta]}=q_1^{c_1}\cdots q_n^{c_n}$，$c_a=\langle\beta,\eta_a\rangle$。并且
\begin{equation}\tag{3.20}
N_{[\beta]}(\eta_a,\eta_b,\eta_c)=n_{[\beta]}\langle\beta,\eta_a\rangle\langle\beta,\eta_b\rangle\langle\beta,\eta_c\rangle\,,
\end{equation}
它们是有理数，称为 Gromov--Witten 不变量。由 (3.19) 可见，三点函数（汤川耦合）的几何解释包含两部分：由三重相交给出的经典部分，以及由瞬子求和给出的量子部分。
